% ============================================================
%  Paper 14: The f(R) = R^(67/48) Gravitational Action and Horndeski EFT Parameters
%  Target: JCAP Compilation (SDGFT Series)
% ============================================================
\documentclass[a4paper,11pt]{article}
\usepackage{jcappub}

\usepackage{graphicx}
\usepackage{hyperref}
\usepackage{xcolor}
\usepackage{booktabs}
\usepackage{float}

% ── SDGFT shared macros ──────────────────────────────────────
\usepackage{sdgft-macros}

% ── Document ─────────────────────────────────────────────────
\begin{document}

\title{The f(R) = R^(67/48) Gravitational Action and Horndeski EFT Parameters}

\author{David A. Besemer}
\affiliation{Independent Researcher}
\emailAdd{david.besemer@sdgft.org}

\abstract{
We formulate the gravitational action of SDGFT as an f(R) = R^(67/48) modified gravity theory. We derive the four Bellini-Sawicki Horndeski parameters, showing that the Planck mass run-rate is alpha_M = 19/86, the braiding is alpha_B = -19/172, and the tensor speed excess is alpha_T = 0, in compliance with the GW170817 constraint.
}

\arxivnumber{SDGFT-2026-14}

\maketitle

\section{Introduction}
Modified gravity theories often introduce arbitrary functions. In SDGFT, the f(R) action is fixed by the effective spectral dimension.

\section{Theoretical Formulation}
Here we formulate the core mathematical relations governing the system:
\begin{equation}
  S = \frac{1}{16\pi G} \int d^4x\,\sqrt{-g}\,R^{67/48}
\end{equation}
\begin{equation}
  \aM = \frac{n-1}{2n-1} = \frac{19}{86} \approx 0.2209, \quad \aB = -\frac{\aM}{2}, \quad \aT = 0
\end{equation}
\section{Horndeski Parameterization}
We compute the Bellini-Sawicki parameters for f(R) = R^n and show that the gravitational wave propagation speed is exactly equal to the speed of light.

\section{Conclusion}
The f(R) = R^(67/48) theory is stable and matches all solar system and gravitational wave constraints.



\begin{thebibliography}{99}
\bibitem{Goldstein_Paper01} D. A. Besemer, ``Axiomatic Foundations of SDGFT,'' SDGFT-2026-01 (in preparation).
\bibitem{Goldstein_Paper04} D. A. Besemer, ``Effective Spectral Dimension and the Banach Fixed-Point Theorem in SDGFT,'' SDGFT-2026-04.
\bibitem{Goldstein_Code} D. A. Besemer, \texttt{sdgft} Python package, \url{https://github.com/david-besemer/sdgft} (2026).
\end{thebibliography}

\end{document}
